\begin{lemma}[Comparison of $\gamma_n$ with Lebesgue measure on a ball, and the null-set equivalence]
Fix $n \in \mathbb{N}_0$, let $\lambda$ denote Lebesgue measure on
$\mathbb{R}^n$ and let $\gamma_n$ be the Gaussian measure of
\noderef{GaussMeasure}.  For $R \in \mathbb{R}$ write
$a_R = \pi^{-n/2}e^{-R^2}$.  Then:
\begin{enumerate}
\item for \emph{every} set $S \subseteq \mathbb{R}^n$, measurable or not,
\[
  \gamma_n(S) \leq \pi^{-n/2}\,\lambda(S) ;
\]
\item for every $R \in \mathbb{R}$ and every measurable
$E \subseteq \overline B(0,R)$,
\[
  a_R\,\lambda(E) \leq \gamma_n(E) ;
\]
\item for every $R, \rho \in \mathbb{R}$ and every $x \in \mathbb{R}^n$ with
$|x| \leq R$ and $0 < \rho \leq 2$,
\[
  a_{R+2}\,\lambda\bigl(\overline B(x,\rho)\bigr)
    \leq \gamma_n\bigl(\overline B(x,\rho)\bigr) ;
\]
\item for every set $A \subseteq \mathbb{R}^n$, measurable or not,
\[
  \gamma_n(A) = 0 \iff \lambda(A) = 0 .
\]
\end{enumerate}
Here $\overline B(x,\rho)$ is the closed ball, and both $\gamma_n$ and $\lambda$
are applied to arbitrary sets as outer measures, i.e.\ through
$\mu(S) = \inf\{\mu(F) : S \subseteq F,\ F \text{ measurable}\}$.  All products
and comparisons are read in $[0,+\infty]$.

\emph{What this adds to \noderef{GaussMeasure}, which only says what
$\gamma_n$ is.}  Clause (1) is a global one-sided comparison, valid for all
sets, and it is the only direction that is globally true.  Clause (2) is the
reverse comparison, which is \emph{false} globally --- the ratio
$\gamma_n(E)/\lambda(E)$ tends to $0$ as $E$ escapes to infinity --- and
becomes true, with a quantitative constant, exactly because $E$ is confined to
a ball; it is this localization that removes any need for covering theory
adapted to the non-doubling measure $\gamma_n$, by reducing a local Gaussian
averaging argument to a Lebesgue one.  Clause (3) is the form in which a
covering argument uses clause (2): the ball $\overline B(x,\rho)$ is not
centred at the origin, and the enlargement from $R$ to $R+2$ is exactly what a
cap of $2$ on the radius of the balls costs.  Clause
(4) is a statement about \emph{all} sets and is not a comparison of sizes at
all: it says the two measures have the same null sets, so that
``$\lambda$-almost everywhere'' and ``$\gamma_n$-almost everywhere'' are
interchangeable, and so that an almost-everywhere identity between continuous
functions for one measure is one for the other.  None of the four is the
density formula, and none follows from it without the argument below.

\emph{The constants are explicit and dimension-dependent.}  Both $\pi^{-n/2}$
and $a_R = \pi^{-n/2}e^{-R^2}$ depend on $n$, and $a_R$ degrades as $R$ grows;
no uniformity in $n$ or in $R$ is claimed.
\end{lemma}
\begin{proof}
By \noderef{GaussMeasure}, $\gamma_n$ is Lebesgue measure weighted by the
density
\[
  \varrho(y) = \pi^{-n/2}e^{-|y|^2} ,
\]
that is, $\gamma_n(E) = \int_E\varrho\,d\lambda$ for every measurable
$E \subseteq \mathbb{R}^n$, the integral being the lower Lebesgue integral of a
nonnegative function.  Note $0 < \varrho(y) \leq \pi^{-n/2}$ for every $y$,
since $0 < e^{-|y|^2} \leq 1$, and $\varrho(y) \geq \pi^{-n/2}e^{-R^2} = a_R$
whenever $|y| \leq R$, since $t \mapsto e^{-t}$ is decreasing and
$|y|^2 \leq R^2$ in that case.

\emph{Step 1: clause (1) for measurable sets.}
Let $E$ be measurable.  By monotonicity of the integral in the integrand,
\[
  \gamma_n(E) = \int_E\varrho\,d\lambda
    \leq \int_E\pi^{-n/2}\,d\lambda
    = \pi^{-n/2}\lambda(E) .
\]

\emph{Step 2: clause (1) for arbitrary sets.}
Let $S \subseteq \mathbb{R}^n$ be arbitrary and let $F$ be any measurable set
with $S \subseteq F$.  By monotonicity of $\gamma_n$ and Step 1,
\[
  \gamma_n(S) \leq \gamma_n(F) \leq \pi^{-n/2}\lambda(F) .
\]
Taking the infimum over all such $F$ --- the family is nonempty, since
$F = \mathbb{R}^n$ qualifies --- and using that multiplication by the finite
constant $\pi^{-n/2}$ commutes with infima of nonempty families in
$[0,+\infty]$, we get
\[
  \gamma_n(S) \leq \pi^{-n/2}\inf\{\lambda(F) : S \subseteq F,\ F \text{
  measurable}\} = \pi^{-n/2}\lambda(S) ,
\]
the last equality being the displayed description of $\lambda$ on arbitrary
sets.  This is clause (1).

\emph{Step 3: clause (2).}
Let $R \in \mathbb{R}$ and let $E$ be measurable with
$E \subseteq \overline B(0,R)$.  If $R < 0$ then $\overline B(0,R) = \emptyset$,
so $E = \emptyset$ and both sides vanish.  Assume $R \geq 0$.  Every $y \in E$
satisfies $|y| \leq R$, hence $\varrho(y) \geq a_R$ there, and therefore
\[
  \gamma_n(E) = \int_E\varrho\,d\lambda
    \geq \int_E a_R\,d\lambda
    = a_R\,\lambda(E) ,
\]
again by monotonicity of the integral in the integrand.  This is clause (2).

\emph{Step 4: clause (3).}
Let $x$ with $|x| \leq R$ and $0 < \rho \leq 2$.  For every
$y \in \overline B(x,\rho)$ the triangle inequality gives
\[
  |y| \leq |x| + |y - x| \leq R + \rho \leq R + 2 ,
\]
so $\overline B(x,\rho) \subseteq \overline B(0,R+2)$.  The closed ball is a
closed, hence measurable, set.  Clause (2) applied with $R+2$ in place of $R$
and $E = \overline B(x,\rho)$ gives
$a_{R+2}\lambda(\overline B(x,\rho)) \leq \gamma_n(\overline B(x,\rho))$,
which is clause (3).

\emph{Step 5: clause (4), the direction $\lambda(A) = 0 \Rightarrow
\gamma_n(A) = 0$.}
Immediate from clause (1): $\gamma_n(A) \leq \pi^{-n/2}\cdot 0 = 0$.

\emph{Step 6: clause (4), the direction $\gamma_n(A) = 0 \Rightarrow
\lambda(A) = 0$.}
Suppose $\gamma_n(A) = 0$.  Since $\gamma_n(A)$ is the infimum of $\gamma_n(F)$
over measurable $F \supseteq A$, there is a measurable $F$ with $A \subseteq F$
and $\gamma_n(F) = 0$.  (Choose measurable $F_k \supseteq A$ with
$\gamma_n(F_k) \leq 1/k$ and set $F = \bigcap_kF_k$, which is measurable,
contains $A$, and has $\gamma_n(F) \leq \gamma_n(F_k) \leq 1/k$ for every $k$,
hence $\gamma_n(F) = 0$.)

Fix $m \in \mathbb{N}$ and put $F_m' = F \cap \overline B(0,m)$, a measurable
subset of $\overline B(0,m)$.  By clause (2) with $R = m$, and by monotonicity
of $\gamma_n$,
\[
  a_m\,\lambda(F_m') \leq \gamma_n(F_m') \leq \gamma_n(F) = 0 .
\]
Since $a_m = \pi^{-n/2}e^{-m^2} > 0$, this forces $\lambda(F_m') = 0$.  Now
$\mathbb{R}^n = \bigcup_{m\in\mathbb{N}}\overline B(0,m)$, so
$F = \bigcup_{m\in\mathbb{N}}F_m'$, and countable subadditivity of $\lambda$
(shared context) gives $\lambda(F) \leq \sum_m\lambda(F_m') = 0$.  Finally
$\lambda(A) \leq \lambda(F) = 0$ by monotonicity.  This completes clause (4),
and the proof.
\end{proof}
